\section{本章小结}
\label{sec:lossless-summary}

本章从“常见符号用短码能省多少”出发，回答了三个层次的问题，每个层次对应各自的约束。

第一层是可解码性。码字互不相同还不够，拼接以后必须能唯一地切开。前缀码让解码器在码树上即时断开码字；唯一可译码允许码字互为前缀，但要付出解码延迟。Kraft–McMillan不等式说明两者能实现的码长集合相同，都由$\sum_x2^{-\ell(x)}\leq1$刻画；这一条件约束的是码长，码字仍须按保证前缀性的方法分配。

第二层是平均码长的最优性。对至少两个符号有正概率的信源，唯一可译码的平均码长不小于熵，等号要求每个码长恰为自信息；按自信息取整的Shannon码平均码长小于熵加一。平均码长与熵之差拆成两部分：码长暗含的概率与真实概率的KL散度，以及Kraft不等式未取等号所浪费的空间。冗余是加权平均，个别符号的码长可以短于自信息。确定源只有在消息长度已知时才能零码长。

第三层是可执行的构造与渐近极限。Huffman编码逐次合并概率最小的两个符号，对给定有限符号表求出平均码长最短的前缀码，贯穿例中为$1.9$位，熵为$1.846$位。

算术编码把整条消息对应到长度等于其概率的区间，码长不超过自信息加两位，需要约定终止规则，两端须使用完全相同的模型与取整方式，模型给实际符号零概率时无法编码。

把符号组成块之后有两种资源制度：零错误的变长块编码使每符号冗余小于$1/n$，代价是最坏码长和码本随块长增长；允许错误的定长块编码借助渐近均分性和典型集，码率高于熵时错误概率趋于零，低于熵时趋于一，失败来自被放弃的小概率块。

在有限资源下，块长带来的收益很快被码本与计算代价抵消；模型失配使平均码长增加KL散度，两端不一致则直接造成解码错误。第~\ref{chap:quantization}章不再要求精确还原，允许重构值与原值有偏差，研究把连续取值装进有限个格子的量化器，并用失真来衡量这种偏差；量化后的格子编号仍可用本章的方法无损压缩。
